\documentclass[11pt,fontset=fandol]{ctexart}

\usepackage[a4paper,margin=2.35cm]{geometry}
\usepackage{enumitem}
\usepackage{hyperref}
\usepackage{xurl}

\hypersetup{
  colorlinks=true,
  linkcolor=black,
  citecolor=black,
  urlcolor=blue
}

\setlist[itemize]{leftmargin=1.8em,itemsep=0.2em,topsep=0.2em}
\Urlmuskip=0mu plus 2mu
\emergencystretch=2em

\title{\textbf{SciTrace-RL：科学智能体执行轨迹、验证与后训练数据层}}
\author{追光计划项目提议}
\date{}

\begin{document}
\maketitle

\section*{1. 具体目标}

构建面向 AI4S Infra 的科学智能体执行账本，解决科学任务中“过程不可追踪、结论不可验证、失败经验不可复用”的问题。系统记录任务目标、工具调用、证据来源、产物哈希、验证结果和失败类型，产出 trace schema、验证器、15-case 评测集、Demo 报告、trace-to-reward 样本、后训练数据包、升级验证包，以及对齐 W3C PROV、Workflow Run RO-Crate 和 OpenTelemetry span 的 provenance bundle。

\section*{2. 核心价值}

Bohrium+SciMaster 指出，规模化科学工作流需要可观察、可复现、可治理的执行基础设施\cite{bohrium_scimaster}；CORE-Bench 和 MLR-Bench 也显示，研究智能体的主要风险是复现失败、虚构实验和无效结论\cite{corebench,mlrbench}。SciTrace-RL 的差异是反证优先：先验证系统能否识别伪造引用、错误机理、虚构计算和过度结论，再输出推荐。

大模型部分聚焦后训练：将工具调用轨迹、验证分数和失败类型转化为 SFT、DPO/GRPO、Process Reward Model、credit assignment 和 tool-router 数据，使科学智能体从真实执行反馈中改进规划、检索、工具调用和自我校验能力。

\section*{3. 可行性简析}

Demo 以锂电电解液添加剂筛选为例，已实现本地原型、DeepSeek 语义审查、15 个评测案例、trace-to-reward、post-training bundle、provenance bundle 和 escalation packet。当前 80\% 案例可由规则验证器和大模型裁判自动处理；20\% 涉及高电压电芯选择、湿实验协议和多物理 tradeoff，明确升级到 Bohrium/Lebesgue 高保真计算、Uni-Lab-OS 实验或专家判断\cite{dp_about,unilabos}。

\begin{thebibliography}{9}

\bibitem{bohrium_scimaster}
Linfeng Zhang et al.
\textit{Bohrium + SciMaster: Building the Infrastructure and Ecosystem for Agentic Science at Scale}.
arXiv:2512.20469, 2025.
\url{https://arxiv.org/abs/2512.20469}

\bibitem{corebench}
Zachary S. Siegel et al.
\textit{CORE-Bench: Fostering the Credibility of Published Research Through a Computational Reproducibility Agent Benchmark}.
arXiv:2409.11363, 2024.
\url{https://arxiv.org/abs/2409.11363}

\bibitem{mlrbench}
Hui Chen et al.
\textit{MLR-Bench: Evaluating AI Agents on Open-Ended Machine Learning Research}.
arXiv:2505.19955, 2025.
\url{https://arxiv.org/abs/2505.19955}

\bibitem{dp_about}
DP Technology.
\textit{Particle Universe / Company and Product Overview}.
\url{https://www.dp.tech/en/about}

\bibitem{unilabos}
DP Technology.
\textit{Uni-Lab-OS}.
\url{https://github.com/dptech-corp/Uni-Lab-OS}

\end{thebibliography}

\end{document}
